\section{本章纲要：一个把富庶变成边防的国家}

北宋不是从一条单线的“盛世”走向一次突然的覆亡。开封的宫城、汴河、仓场与禁军，把五代以来最容易失控的军权重新系到中枢；四川的纸币、江南的粮食、河北的河工和西北的堡寨，又把这个中枢系到彼此条件全然不同的地方。它因此能在九百六十至九百七十九年的战争中接收大部分旧政权，却不能把燕云、河东以北和西北边地当作等待填色的空白。

本章把宋放回与辽、西夏、金同时行动的世界。澶渊的银绢、庆历的和议和熙河的堡寨，并非“对外”章节的插曲；它们决定了京师须筹多少钱粮、地方须承受何种征发、官僚又以什么语言争论富国强兵。新法也不只是王安石与反对者的立场之争：均输、青苗、免役和市易试图使国家在运输、信用和劳役市场中取得更稳定的接口，随之而来的则是执行者、记录与地方承受力的问题。

读者会先看建国、统一与宋辽关系如何造出昂贵的北方国家，再看改革、城市和区域市场怎样扩展并暴露这种国家能力，最后回到海上之盟与靖康。那些事件不能预先被写成必然的“宋亡原因”。它们是几套制度在战争、外交、河流、市场和继承政治中相互牵动的结果。正史本纪与《续资治通鉴长编》给出朝廷年序，《宋会要辑稿》保存机构与条令，地方遗址、文书和钱币则使城市、流通和区域差异较为可见；没有任何一种材料足以单独替所有人说话。
